% SPDX-FileCopyrightText: 2026 Will Cook
% SPDX-License-Identifier: CC-BY-4.0
\documentclass[11pt]{article}
\usepackage[a4paper,margin=28mm]{geometry}
\usepackage{amsmath,amssymb,amsthm,mathtools,microtype,booktabs}
\usepackage[colorlinks=true,allcolors=blue]{hyperref}
\usepackage{xurl}
\newtheorem{theorem}{Theorem}
\newtheorem{lemma}[theorem]{Lemma}
\newtheorem{corollary}[theorem]{Corollary}
\newtheorem{proposition}[theorem]{Proposition}
\newtheorem{remark}[theorem]{Remark}
\newcommand{\N}{\mathbb N}
\newcommand{\Z}{\mathbb Z}
\newcommand{\R}{\mathbb R}
% Immutable formal-source revision for the Lean ingredients cited below.
\newcommand{\commit}{8d6596fbde2c4aacf946adec4b0226ba1a97545d}
\newcommand{\commitshort}{8d6596fbde2c}
\newcommand{\repobase}{https://github.com/wcook04/plectis-erdos/blob/\commit}
% Same first-page production disclosure as the synthesis papers.
\newcommand{\noteprovenance}{%
  \begingroup\renewcommand{\thefootnote}{*}%
  \footnotetext{\emph{Authorship and AI use.}
Will Cook built and directed the research infrastructure and reviewed the claims
when he could. AI agents did most of the research and drafting. Cook did not
independently verify every claim.}\endgroup}
\title{An Exact Capacity Criterion for Series\\with Eventual Congruences}
\author{Will Cook\textsuperscript{*}}
\date{20 September 2026}
\begin{document}
\maketitle
\noteprovenance
\begin{abstract}
Let $Q_n$ be positive integers with $Q_n\mid Q_{n+1}$ and
$Q_{n+1}\ge2Q_n$, and let $F_n\ge0$ be integer allowances with
$\sum F_n/Q_n<\infty$. Series whose nonnegative integer digits satisfy these
bounds eventually and are eventually divisible by every fixed integer
represent an interval exactly when $Q_N\sum_{n>N}F_n/Q_n\to\infty$.
In that case every nonnegative real is represented; otherwise the set of
values is null and meagre. The sufficient direction
also imposes eventual divisibility of the cumulative digit sums, with
cutoffs common to every target in the interval. Arbitrary oscillation,
zero allowances and unbounded denominator ratios are permitted. For
factorial denominators and allowances $n^c$ on a fixed support, the criterion
gives the sharp gap threshold $\text{gap}<c$; its contrast with the
unrestricted threshold $\text{gap}\le c$ occurs exactly at integer exponents.
The proof combines common continuation intervals for every cumulative residue
with the arithmetic lattice of prefixes. A separate common-divisor criterion
isolates an obstruction to rationality for bounded-carry Cantor series.
\end{abstract}

\section{The exact criterion}
A factorial expansion usually permits a digit of size about $n$ at position
$n$. If more digits are allowed, can most positions be discarded? What changes
if each fixed integer must eventually divide every digit?

Let $Q_n$ be positive integers, $n\ge1$, such that $Q_n\mid Q_{n+1}$ and
$Q_{n+1}\ge2Q_n$. Let $F_n\in\Z_{\ge0}$ satisfy
$\sum_nF_n/Q_n<\infty$, and put
\[
 U_N=Q_N\sum_{n>N}\frac{F_n}{Q_n}.
 \tag{C}\label{capacity:eq:capacity}
\]
Define $E_Q(F)$ to be the set of sums $\sum_{n\ge1}e_n/Q_n$, allowing
nonnegative integer digits $e_n\le F_n$ eventually and requiring, for each
fixed $q\ge1$, that $q\mid e_n$ eventually. The cutoffs and finite initial
digits may depend on the represented value.

\begin{theorem}\label{capacity:thm:capacity}
If $U_N\to\infty$, then $E_Q(F)=[0,\infty)$; otherwise $E_Q(F)$ is null
and meagre. In particular interval filling is equivalent to $U_N\to\infty$.
In the positive case, every nonnegative target may be represented with
$e_n=0$ before any prescribed cutoff, at most one exception to $e_n\le F_n$,
and
\[
 q\mid e_n,\qquad q\mid\sum_{k<n}e_k
 \quad\text{for all sufficiently large }n,
\]
with congruence cutoffs depending on $q$ but independent of the target.
If the allowance must hold at every position, a nondegenerate interval is
still represented with both congruences and a common cutoff for each $q$.
\end{theorem}
The criterion compares \emph{all} future capacity with the prefix lattice.
It requires no monotonicity, regular variation or gap bound on the allowances,
and no upper bound on the denominator ratios. It does not require the $Q_n$
to clear every rational denominator. Section~\ref{capacity:sec:capacity-proof} proves
the theorem using the same residue-feedback mechanism already present in
the public Lean repository. The capacity estimates and complete criterion
are ordinary proofs; the abstract feedback endpoint is kernel-checked
separately. Both geometric and factorial denominators satisfy the hypotheses.

\subsection{The sharp support threshold}
Let $S\subseteq\N_{\!>0}$ and $c>0$. Define $E(S,c)$ to be the set of sums
\[
 x=\sum_{n\ge1}\frac{e_n}{n!},\qquad e_n\in\Z_{\ge0},\quad
 e_n=0\ (n\notin S),\quad e_n\le n^c\text{ eventually},
 \tag{1}\label{capacity:eq:class}
\]
subject to
\[
 \text{for every }q\ge1,\qquad q\mid e_n\text{ for all sufficiently large }n.
 \tag{2}\label{capacity:eq:congruence}
\]
The cutoffs and the finite initial digits in this definition may depend on $x$.
All series converge, since finitely many unrestricted digits do not affect
convergence. A permitted position need not carry a nonzero digit.

\begin{theorem}\label{capacity:thm:main}
If $S$ is finite, $E(S,c)$ is countable. If $S$ is infinite, enumerate it as
$n_0<n_1<\cdots$.
\begin{enumerate}
\item For $0<c\le1$, $E(S,c)$ is null and meagre.
\item For $c>1$, $E(S,c)$ contains a nondegenerate interval if and only if
\[
 n_j-n_{j-1}<c\qquad\text{for all sufficiently large }j.
 \tag{3}\label{capacity:eq:gaps}
\]
If this condition fails, $E(S,c)$ is null and meagre.
\end{enumerate}
When \eqref{capacity:eq:gaps} holds, one may require $e_n=0$ before any prescribed
cutoff and impose both
\[
 q\mid e_n,\qquad q\mid\sum_{k<n}e_k
 \tag{4}\label{capacity:eq:both}
\]
for all sufficiently large $n$, with cutoffs depending on $q$ but independent
of $x$ throughout the constructed interval.
\end{theorem}
The theorem is an ordinary mathematical proof. The associated formal sources
check the general digit-feedback construction and, separately, the
common-divisor carry obstruction of Section~\ref{capacity:sec:carry}; they do not
formalise this gap classification or its measure argument. The source and
attribution account is in Section~\ref{capacity:sec:sources}.

For example, with allowance $n^2$, every second position suffices without
congruences. Under \eqref{capacity:eq:congruence}, infinitely many omitted positions
already force a null set. With allowance $n^{2.01}$, every second position
again suffices, even under \eqref{capacity:eq:both}. Thus the strict inequality in
\eqref{capacity:eq:gaps} is essential.

\begin{corollary}\label{capacity:cor:density}
For $c>1$, the least possible asymptotic density of a fixed permitted support
that fills an interval under \eqref{capacity:eq:congruence} is
$1/(\lceil c\rceil-1)$. Every such support has lower density at least this
value, and an arithmetic progression attains it. A density bound alone is
not sufficient: even rare gaps of length $\lceil c\rceil$ prevent interval
filling if they occur infinitely often.
\end{corollary}
\begin{proof}
Put $d=\lceil c\rceil-1$. Eventual gaps at most $d$ give
$\liminf_{N\to\infty}|S\cap[1,N]|/N\ge1/d$. An arithmetic progression of
step $d$ satisfies the theorem.
\end{proof}

\section{Proof of the capacity criterion}\label{capacity:sec:capacity-proof}
The necessity and sufficiency use different parts of the expansion. Prefixes
supply an arithmetic obstruction; a common interval of possible future
remainders makes the residue choices compatible with exact representation.

\subsection{Bounded capacity along a subsequence}
If $U_N$ does not tend to infinity, there are a constant $B>0$ and arbitrarily
large $N$ with $U_N\le B$. Choose one integer $q>B$. Fix a cutoff $K$ after
which the allowance and divisibility by $q$ hold, and fix the finite prefix,
of value $y_0$. The set $V$ of resulting sums is compact by summability and
the finite choices at each remaining position. For every such $N>K$,
\[
 V\subseteq y_0+\frac q{Q_N}\Z+[0,B/Q_N].
\]
Indeed the scaled prefix after $K$ is a multiple of $q$ and the remaining
tail is at most $U_N/Q_N$. These grids leave holes inside every fixed open
interval for all sufficiently large selected $N$, so $V$ is nowhere dense.
For a bounded interval $J$ they also give
$|V\cap J|\le(B/q)|J|+2B/Q_N$.
Passing to the subsequence and covering $V$ by finitely many intervals with
total length at most $|V|+\varepsilon$ gives
$|V|\le(B/q)(|V|+\varepsilon)$, hence $|V|=0$.
A countable union over the cutoffs and integer prefixes covers $E_Q(F)$.

\subsection{One continuation interval for every cumulative residue}
Suppose $U_N\to\infty$ and put $h_N=\inf_{k\ge N}U_k$. For sufficiently
large $n$, let $M_n$ be the largest factorial at most
$\min\{n,\sqrt{h_{\lfloor n/2\rfloor}}\}$, and set the finitely many
earlier moduli equal to one. These positive integers are nested under
divisibility, tend to infinity, and are eventually divisible by each fixed
integer. The delayed lower envelope controls the whole future cost of
rounding, even when the denominator ratios are unbounded. Indeed,
$Q_N/Q_k\le2^{N-k}$, so for sufficiently large $N$,
\[
 \begin{aligned}
 C_N:=Q_N\sum_{k>N}\frac{M_k}{Q_k}
 &\le\sqrt{h_N}\sum_{k=N+1}^{2N}2^{N-k}
       +\sum_{k>2N}k\,2^{N-k}\\
 &\le\sqrt{h_N}+(2N+2)2^{-N}.
 \end{aligned}
\]
The first bound uses $\lfloor k/2\rfloor\le N$ when $k\le2N$; the second
uses $M_k\le k$. Since $U_N\ge h_N\to\infty$ and
$M_N\le\sqrt{h_N}$, this proves $U_N-2C_N\ge M_N$ for every sufficiently
large $N$. Start beyond this point and any prescribed cutoff.

Keep only the active positions $S=\{n:F_n\ge2M_n\}$ after this cutoff.
Define two common tail bounds
\[
 \alpha_N=\sum_{\substack{n>N\\n\in S}}\frac{M_n}{Q_n},\qquad
 \beta_N=\sum_{\substack{n>N\\n\in S}}\frac{F_n-M_n}{Q_n}.
\]
Both are nonnegative and tend to zero. Discarding an inactive position loses
less than $2M_n$; reserving both margins at an active position loses exactly
$2M_n$. Thus the preceding estimate gives the decisive overlap inequality
\[
 Q_N(\beta_N-\alpha_N)\ge U_N-2C_N\ge M_N.
 \tag{O}\label{capacity:eq:overlap}
\]
In particular the active set is infinite and each continuation interval
has positive width.

Choose a target in $[\alpha_{N_0},\beta_{N_0}]$ and set all preceding digits
to zero. Suppose the current unweighted cumulative sum is $C$ and the
remaining target before position $n$ lies in
$[\alpha_{n-1},\beta_{n-1}]$. At an inactive position choose zero.
At an active position, the permitted digits in the residue class
$e\equiv-C\pmod{M_n}$, $0\le e\le F_n$, form a nonempty arithmetic
progression of step $M_n$. Its least digit is less than $M_n$ and its
greatest digit exceeds $F_n-M_n$. By (O), the intervals
\[
 \frac e{Q_n}+[\alpha_n,\beta_n]
\]
for these digits overlap. Their union contains
\[
 \left[\frac{M_n}{Q_n}+\alpha_n,
       \frac{F_n-M_n}{Q_n}+\beta_n\right]
   =[\alpha_{n-1},\beta_{n-1}].
\]
Choose a digit whose interval contains the remaining target. The next
remainder stays between $\alpha_n$ and $\beta_n$, and these bounds tend to
zero, proving exact representation. This is precisely the common-interval
feedback construction; no target-independent ordinary block sum is required.

After an active position $n$, the cumulative sum is divisible by $M_n$.
Each later active digit is divisible by the previous active modulus,
because both adjacent cumulative sums are; the moduli are nested.
For a fixed $q$, take one active position with $q\mid M_n$. After that
position, all individual digits and all preceding cumulative sums are
divisible by $q$. This cutoff is common to all targets in the interval.

\subsection{From one interval to every nonnegative target}
The finite-exception convention in $E_Q(F)$ strengthens the conclusion.
Given $y>0$, choose $N$ arbitrarily late with $\alpha_N<y$ and put
\[
 p=M_N\left\lfloor\frac{Q_N(y-\alpha_N)}{M_N}\right\rfloor.
\]
Then $p\ge0$, $M_N\mid p$, and (O) gives
$y-p/Q_N\in[\alpha_N,\beta_N]$. Set $e_N=p$ and all earlier digits to
zero. Run the same continuation construction after $N$, with initial
cumulative sum $C=p$; its interval covering worked for every $C$.
All later digits obey their allowances. The zero target uses zero digits.

The congruence cutoffs can also be common to all $y$. For fixed $q$, choose
an active position $j$ with $q\mid M_j$. If the target's exceptional index
$N$ precedes $j$, the repair at $j$ gives the required divisibilities after
$j$. If $N\ge j$, all preceding digits are zero, $q\mid M_N\mid p$, and
all subsequent repairs preserve divisibility by $q$. Thus $j+1$ is a valid
cutoff for every target. The allowance-exception index itself may depend
on the target. This completes Theorem~\ref{capacity:thm:capacity}.

The growth hypothesis has content. If repeated denominators are allowed,
take $Q_n=2^{\lfloor\sqrt n\rfloor}$ and $F_n=1$. The allowance series
converges, while $U_N\ge\lfloor\sqrt N\rfloor$ by counting the next complete
denominator block. Yet eventual divisibility by two forces all sufficiently
late digits to vanish, so the attainable set is countable. Imposing all
cumulative congruences as well leaves only zero. Thus nestedness alone does
not justify the criterion.

\section{Why a large gap prevents interval filling}\label{capacity:sec:negative}
The decisive obstruction is local. At a position $N$, the whole preceding
factorial sum lies on a lattice of spacing $1/N!$. Eventual divisibility
widens that spacing to $q/N!$, up to a fixed translation. After a sufficiently
long gap, all remaining permitted digits reach only a bounded multiple of
$1/N!$. Choosing one fixed $q$ larger than that bound leaves holes at
arbitrarily small scales.

Here are the details, including the target-dependent cutoffs in (1)--(2).
Put $r=\lceil c\rceil$. Suppose infinitely many successive $S$-gaps have
length at least $r$, and let $N$ run through the positions immediately before
these gaps. For large $k$,
\[
 \frac{(k+1)^c/(k+1)!}{k^c/k!}
 =\frac{(1+1/k)^c}{k+1}\le\frac12.
\]
Consequently, for sufficiently large such $N$,
\begin{align}
 N!\sum_{\substack{k>N\\k\in S}}\frac{k^c}{k!}
 &\le 2N!\frac{(N+r)^c}{(N+r)!}\\
 &\le 2^{c+1}N^{c-r}\le B,\qquad B=2^{c+1}.
 \tag{5}\label{capacity:eq:tailbound}
\end{align}
In the second inequality we used $N\ge r$ and
$(N+1)\cdots(N+r)\ge N^r$.

The bounded capacity along this subsequence invokes the necessity argument
of Theorem~\ref{capacity:thm:capacity}, with allowances
$F_n=\lfloor n^c\rfloor$ on $S$ and zero elsewhere. It gives nullity and
meagreness even with target-dependent congruence and allowance cutoffs.

If $c\le1$, every support gap is at least $r=1$, so this proves the first
part of Theorem~\ref{capacity:thm:main}. If $c>1$ and (3) fails, it proves the negative
part. Notice that infrequent large gaps are enough; average digit counts do
not detect this obstruction.

\section{A digit that also corrects the cumulative residue}
The following construction supplies the positive direction. It is useful
beyond factorial weights.

\begin{lemma}\label{capacity:lem:feedback}
Let $w_j>0$, let positive integers $M_j$ satisfy $M_{j-1}\mid M_j$, and let
$F_j\ge0$. Suppose $M_jw_j\to0$ and, for $j\ge1$,
\[
 2M_j\le M_{j-1}\frac{w_{j-1}}{w_j},\qquad
 2M_{j-1}\frac{w_{j-1}}{w_j}\le F_j.
 \tag{7}\label{capacity:eq:step}
\]
Every $y\in[M_0w_0,2M_0w_0]$ has a representation
$y=\sum_{j\ge1}b_jw_j$ with integers $0\le b_j\le F_j$ such that
\[
 M_j\mid\sum_{i\le j}b_i,\qquad M_{j-1}\mid b_j.
\]
\end{lemma}
\begin{proof}
Start with residual $R_0=y$ and cumulative sum $C_0=0$. Suppose
$M_{j-1}w_{j-1}\le R_{j-1}\le2M_{j-1}w_{j-1}$ and
$M_{j-1}\mid C_{j-1}$. Put $u=R_{j-1}/w_j$ and let
$v$ be the least nonnegative residue of $-C_{j-1}$ modulo $M_j$. Choose
\[
 b_j=v+M_j\left\lfloor\frac{u-M_j-v}{M_j}\right\rfloor.
 \tag{8}\label{capacity:eq:selector}
\]
Since $u\ge2M_j$ and $0\le v<M_j$, this integer is nonnegative.
The floor identity gives $M_j\le u-b_j<2M_j$; also $b_j\le u\le F_j$.
Thus $R_j=R_{j-1}-b_jw_j$ lies in
$[M_jw_j,2M_jw_j)$, while $C_j=C_{j-1}+b_j$ is divisible by $M_j$.
The incoming modulus divides both $C_{j-1}$ and $M_j$, hence also $b_j$.
Finally $R_j\to0$, so the partial sums converge to $y$.
\end{proof}
The residue in (8) may depend on $y$. Only the schedule of moduli needs to
be common to all targets. Requiring target-independent ordinary block sums
would impose an unnecessary restriction on the construction.

\section{The positive direction and the arithmetic cost}
Suppose $c>1$ and the support gaps are eventually at most
$d=\lceil c\rceil-1$. Discard a finite prefix and write
\[
 r_j=\frac{n_j!}{n_{j-1}!},\qquad F_j=n_j^c,\qquad w_j=\frac1{n_j!}.
\]
Then $r_j\ge n_j\to\infty$, while $r_j\le n_j^d$ and therefore
$F_j/r_j\ge n_j^{c-d}\to\infty$. Define
\[
 h_j=\inf_{k\ge j}\min\left\{\frac{\sqrt{r_k}}2,
                                  \frac{F_k}{2r_k}\right\}.
\]
This lower envelope is nondecreasing and tends to infinity. Begin sufficiently
late that $h_1\ge1$, set $M_0=1$, and let $M_j$ be the largest factorial at
most $h_j$. These moduli are nested and eventually divisible by every fixed
integer. Moreover,
\[
 2M_j\le\sqrt{r_j}\le r_jM_{j-1},\qquad
 2r_jM_{j-1}\le F_j.
\]
The second inequality uses $M_{j-1}\le h_j$, also valid for $j=1$.
Finally,
$M_j/n_j!\le\sqrt{r_j}/(2n_j!)\le1/(2\sqrt{n_j!})\to0$.
Lemma~\ref{capacity:lem:feedback} fills $[1/n_0!,2/n_0!]$ using the positions $n_j$,
$j\ge1$. Put zero digits elsewhere. Once $q\mid M_J$, all later digits and
cumulative sums at the original integer indices have the divisibilities in
(4). This proves the positive direction, including common cutoffs.

For comparison, remove condition (2) and call the resulting attainable set
$E_0(S,c)$.
\begin{proposition}\label{capacity:prop:unrestricted}
For an infinite support and $c\ge1$, $E_0(S,c)$ contains an interval exactly
when its successive gaps are eventually at most $\lfloor c\rfloor$.
Otherwise it is null and meagre. For $0<c<1$ it is always null and meagre.
\end{proposition}
\begin{proof}
For necessity repeat Section~\ref{capacity:sec:negative} with
$r=\lfloor c\rfloor+1$. The bound in (5) now tends to zero. With lattice
spacing $1/N!$, it is eventually bounded by $1/(2N!)$, so the same compact-set
argument applies with $q=1$, $B=1/2$.
For sufficiency, $r_j\le n_j^c$ along the support. The ordinary mixed-radix
expansion using digits $0\le b_j<r_j$ fills $[0,1/n_0!]$. To see this
directly, the maximum tail after $n_j$ is $1/n_j!$, since
$(r_k-1)/n_k!=1/n_{k-1}!-1/n_k!$. The adjacent digit intervals therefore
meet, and their lengths tend to zero.
\end{proof}
Thus the difference between the two gap thresholds occurs exactly at integer
$c$. At $c=m\ge2$, imposing eventual congruences raises the least support
density from $1/m$ to $1/(m-1)$. At $c=1$ it destroys interval filling entirely.
With quadratic allowances, deleting only the positions $2^k$ destroys
interval filling, although the permitted support still has density one.
This is the simplest example of why support density loses decisive information.
These statements concern intervals of values on a common permitted support;
they are not lower bounds for representing one specially chosen value.

\section{A common-divisor test for irrationality}\label{capacity:sec:carry}
The distinction between small and large allowances also appears directly in
rationality. Let $b_n\ge2$ be integers, put $Q_0=1$ and
$Q_n=b_1\cdots b_n$, and consider a Cantor series.
\begin{theorem}\label{capacity:thm:gcd}
Suppose $0\le a_n\le A$ and $e_n\ge0$ are integers, $e_n\le Cb_n$
eventually, and $a_n+e_n$ is nonzero infinitely often. If
\[
 \limsup_{n\to\infty}\gcd(b_n,e_n)=\infty,
\]
then $\sum_{n\ge1}(a_n+e_n)/Q_n$ is irrational.
\end{theorem}
\begin{proof}
For large $n$, the scaled tail satisfies
\[
 0<Q_n\sum_{k>n}\frac{a_k+e_k}{Q_k}\le A+2C.
\]
Indeed, the bounded $a_k$ contribute at most
$A\sum_{j\ge1}2^{-j}=A$, and $e_k\le Cb_k$ contributes at most
$C\sum_{j\ge1}2^{-(j-1)}=2C$.
If the total were $p/D$, the numbers
\[
 T_n=DQ_n\left(\frac pD-\sum_{k\le n}\frac{a_k+e_k}{Q_k}\right)
\]
would be positive integers eventually bounded by $D(A+2C)$. No hypothesis
that $D\mid Q_n$ is needed: the extra factor $D$ clears it. The recurrence
\[
 b_nT_{n-1}=D(a_n+e_n)+T_n
\]
implies $\gcd(b_n,e_n)\mid Da_n+T_n$. The integer on the right is positive
and eventually at most $D(2A+2C)$, contradicting the unbounded gcd.
\end{proof}
For factorial denominators, $b_n=n$, and eventual divisibility of $e_n$ by
each fixed integer makes $\gcd(n,e_n)$ arbitrarily large arbitrarily late:
choose a late multiple of that integer. Hence bounded nonnegative digits that
are nonzero infinitely often cannot be rationalised by nonnegative $O(n)$
corrections satisfying (2). Conversely any allowance $F(n)$ with
$F(n)/n\to\infty$ permits interval filling under (4): apply the positive
construction with $r_n=n$, replacing $n^c$ by $F(n)$.

The arithmetic hypothesis cannot be replaced by $b_n\to\infty$, even if both
congruences (4) are retained. Set $a_n=1$, $e_1=2$, and
$e_n=(n+1)!-n!$ for $n\ge2$, and let $b_n=e_n+2$.
Then $\sum_{k\le n}e_k=(n+1)!$, so (4) holds, and $e_n<b_n$. Nevertheless
\[
 \sum_{n\ge1}\frac{a_n+e_n}{Q_n}
 =\sum_{n\ge1}\frac{b_n-1}{Q_n}=1.
\]
Here $\gcd(b_n,e_n)\le2$. This example separates rapid denominator growth
from the arithmetic obstruction used in Theorem~\ref{capacity:thm:gcd}.

\section{Sources, formal correspondence, and further questions}\label{capacity:sec:sources}
The starting point is the sparse perturbation construction accompanying
Erd\H{o}s Problem \#251 in this repository, especially its
\href{../251/erdos-251-prime-gap-dyadic-series.pdf}{short paper} and
\href{\repobase/lean/ErdosProblems/Erdos251/ResidueFeedbackCore.lean}{\texttt{ResidueFeedbackCore.lean}}.
The latter already proves residue-dependent selection and an abstract infinite
sum endpoint. An operator-supplied review supplied the form of
Lemma~\ref{capacity:lem:feedback}, the sharp exponential support constants, and the
linear/superlinear factorial contrast. Those ingredients are credited to that
review, not presented as discoveries of this note. The extensions developed here are the exact capacity criterion on arbitrary
strict integer divisibility chains, the factorial support classification and
its integer-exponent comparison, and the common-divisor formulation.

Airey, Mance and Vandehey already use digit sets eventually divisible by every
fixed integer while retaining asymptotically full digit entropy
\cite[Section 6, p.~1321]{amv2015}. Their theorem concerns normality and
Hausdorff dimension for chosen Cantor bases. Here a fixed divisibility chain and arbitrary summable allowances are given,
and the conclusion distinguishes interval filling from nullity and meagreness. These are elementary arguments in the
classical theory of Cantor series and achievement sets; historical novelty of
the exact classification is not established. Classical interval
covering is background, rather than a contribution claimed here. A literature
comparison and the reproducible checks are recorded in
\href{\repobase/research/experiments/sparse_interpolation/README.md}{the accompanying research record}.
The work was developed with AI assistance and mathematical cross-checking by
separate agent passes; that does not constitute independent expert review.

The formal module
\href{\repobase/lean/ErdosProblems/Synthesis/CongruenceInterpolation.lean}{\texttt{CongruenceInterpolation.lean}}
uses the existing feedback module and states the common-divisor obstruction
for a real carry recurrence. The analytic identification of that recurrence
with the Cantor series, and the capacity and support classifications, are
the ordinary proofs above. The module
\href{\repobase/lean/ErdosProblems/Synthesis/FeedbackContinuation.lean}{\texttt{FeedbackContinuation.lean}}
reuses the existing interval-feedback endpoint and proves eventual individual
and cumulative divisibility from nested cofinal moduli; choosing the moduli
and continuation intervals remains part of the ordinary proof. See the research record for the exact build status and source
revision. No claim about any of the eight Erd\H{o}s programmes changes;
in particular factorial denominators $n!$ here are not $n!-1$ from \#68.

The capacity criterion already covers non-power and oscillating allowances.
For factorial gaps of fixed length $m$, a bounded multiple of $n^m$ still
gives the lattice obstruction, whereas $n^mL(n)$ with $L(n)\to\infty$
permits interval filling. A further question concerns the null case:
what finer tail data determine its Hausdorff dimension? The criterion itself
does not separate dimension zero from full-dimensional null sets. Outside
integer divisibility chains the prefix lattice changes, so no corresponding
necessity is asserted here.
\begin{thebibliography}{9}
\bibitem{amv2015} D.~Airey, B.~Mance and J.~Vandehey,
\emph{Normality preserving operations for Cantor series expansions and
associated fractals. II}, New York J. Math. \textbf{21} (2015), 1311--1326.
\url{https://nyjm.albany.edu/j/2015/21-60v.pdf}.
\end{thebibliography}
\end{document}
